\section{Research Method}
\label{sec:method}
\subsection{Study design and common evaluation}
We organize the research into effect attribution, controlled transfer, and stage diagnosis. The basic comparison unit is a target question with a fixed identifier, choices, and reference answer. Accuracy counts an unextractable answer or a recorded failed response as incorrect. Paired analyses align question identifiers rather than file order. Repeated calls to the same question characterize execution variability; they are not additional independent questions.

\paragraph{Datasets.}
GPQA Main contains 448 graduate-level scientific questions \citep{rein2023gpqa}: the local version has 78 biology, 183 chemistry, and 187 physics items. We use Main, not the 198-item Diamond subset. MuSR supplies narrative reasoning tasks \citep{sprague2024musr}: murder mysteries require combining clues, object placements require tracking locations and agents' knowledge, and team allocation requires satisfying assignment constraints. Unlike GPQA, much of the information needed for MuSR is provided in the narrative. This makes it useful for constructing sources with altered entities but preserved solution relations.

\subsection{Phase I: attributing helpful and harmful effects}
\label{sec:phase1method}
We compare concept-focused teacher analogies with non-analogical explanations, random analogies, and solver-generated reasoning. The teacher content is generated by DeepSeek V4 Flash. The solvers are Llama-3.1-8B-Instruct, Qwen3-8B, Mistral-7B-Instruct-v0.3, and Gemma-4-31B-it. We compare one 300-word analogy with one 600-word analogy, and compare one, two, and three analogies under a nominal total budget of 600 words. Non-analogical explanations have nominal budgets of 300 or 600 words. Random sources are drawn from the same or a different scientific domain. Chain-of-thought (CoT) and self-analogy prompts provide baselines without an external teacher example \citep{wei2022cot,kojima2022zero}.

The controls answer different questions. Length comparisons assess whether more text consistently helps. Explanations assess whether a benefit requires an analogical presentation. Random sources test whether the intended correspondence outperforms unrelated example-shaped content. Multiple examples test the opportunity for comparative abstraction; they do not directly establish that a schema was induced. CoT and self-analogy compare teacher assistance with the solver's own strategies.

For an item-matched baseline $b$ and condition $c$, we decompose the change in accuracy into
\begin{equation}
 p_c-p_b = H_{b\rightarrow c}-D_{b\rightarrow c},
 \label{eq:help-harm}
\end{equation}
where $H$ is the fraction of shared items changing from incorrect to correct, and $D$ is the fraction changing from correct to incorrect. Both use the same shared-item denominator. This retains information that a net improvement discards. We also record explicit textual references to the example, analogy, or teacher. A reference can reject the example; neither its presence nor its absence establishes successful structural use.

Llama, Qwen, and Mistral have 448 retained responses per reported condition. Gemma has 435--447, depending on the condition; missing question records are not fabricated, and paired contrasts use the intersection of available identifiers. Its overall accuracies therefore warrant particular care when comparing conditions with different retained sets.

\subsection{Phase II: testing structural transfer}
\label{sec:phase2method}
For each MuSR target, GPT-5.6-Sol constructs a new source with different entities and context but the intended solution-relevant relational structure. The generator sees the target reference answer to guide source construction; the solver does not receive that answer. Each source contains a narrative, question, choices, and a separate source-answer field. Structural suitability and the absence of an explicit target-answer disclosure are the intended review criteria. Preparation includes format checks, manual review of difficult examples, and a small model-based assessment, rather than exhaustive human validation.

The evaluated set consists of the 103 completed source records available for the main experiment: 61 murder mysteries, 21 object-placement questions, and 21 team-allocation questions. It is not the full MuSR dataset or a documented representative random sample. Random controls reassign the same sources one-to-one among targets with equal choice counts, with no source assigned to its own target. Assignment reduces same-story and same-task matches where possible and considers length differences. Four five-choice object-placement items retain same-story controls because alternatives are unavailable.

Two factors define four conditions: a corresponding or random source, and a provided or withheld source answer. We denote these by $A$ (corresponding, withheld), $B$ (corresponding, provided), $C$ (random, withheld), and $D$ (random, provided). Withholding removes only the source-answer field; the source narrative, question, and choices remain. The target and instruction are fixed. Sources are labeled neutrally as an ``additional example.'' The maximum output is 2,048 tokens and temperature is zero in all four conditions.

We measure the selective source-answer benefit by
\begin{equation}
 \Delta_{\mathrm{transfer}}=(p_B-p_A)-(p_D-p_C).
 \label{eq:interaction}
\end{equation}
A positive interaction means the source answer is more beneficial when attached to an intended corresponding source. We report all four accuracies: an interaction can arise from a benefit in $B$, damage in $D$, or both. A zero interaction does not imply that structural source content has no value.

The four Llama solvers are 3.2-1B, 3.2-3B, 3.1-8B, and 3.1-70B Instruct, each evaluated in three runs. Qwen3-14B and Qwen3-32B provide two-run comparisons outside the Llama family. Run standard deviations describe variability, not confidence intervals. Because random controls permute the same sources, corpus-wide narrative-length distributions are identical (mean 318.99 words), although pairwise and task-specific lengths can differ. We shuffle source choices and update answer indices, then stratify $B$ by whether the source and target correct-answer positions match. This probes index-copying behavior; the strata contain different questions, rather than counterfactual orderings of the same item.

\subsection{Phase III: the stage diagnostic framework}
\label{sec:framework}
The framework uses four fresh model calls. Only the recorded output permitted by each stage contract passes to the next call; conversational state is not carried over. The natural trajectory never receives the target reference answer. Figure~\ref{fig:framework} summarizes the flow, and Table~\ref{tab:contracts} specifies the contracts.

\begin{figure}[t]
\centering
\begin{tikzpicture}[font=\small,>=Stealth,
 stage/.style={draw,rounded corners,align=center,minimum width=6.3cm,minimum height=.65cm},
 node distance=.35cm]
\node[stage] (r) {Retrieval: three solved sources};
\node[stage,below=of r] (m) {Mapping: shared roles and relations};
\node[stage,below=of m] (s) {Schema induction: rule and applicability};
\node[stage,below=of s] (a) {Adaptation: instantiate and solve target};
\draw[->] (r)--(m); \draw[->] (m)--(s); \draw[->] (s)--(a);
\node[above=.18cm of r] {Target enters retrieval};
\node[below=.18cm of a] {Target re-enters adaptation};
\end{tikzpicture}
\caption{The externally controlled four-stage flow. Mapping receives only retrieval output; schema induction receives only mapping output. Adaptation receives the target and schema.}
\label{fig:framework}
\end{figure}

\begin{table*}[t]
\centering
\small
\begin{tabularx}{\textwidth}{@{}p{2.15cm}p{3.8cm}X@{}}
\toprule
Stage & Permitted input & Required output and adequacy criterion \\
\midrule
Retrieval & Target problem and choices & Three solved source problems, including answers and explanations; sources must support a solution-relevant relational structure. \\
Mapping & Recorded retrieval output only & Corresponding source roles and relations, source suitability, and material differences; relationships must be consistent and informative. \\
Schema induction & Recorded mapping output only & An abstract solution rule, roles, relations, applicability conditions, and limits; the schema must preserve enough structure to guide a new solution. \\
Adaptation & Target and necessary schema content & Instantiate schema roles, check target constraints, make necessary adjustments, and give a final answer supported by the target. \\
\bottomrule
\end{tabularx}
\caption{Stage contracts operationalize distinctions from cognitive theory. Mapping is across source examples; the schema is connected to the target during adaptation.}
\label{tab:contracts}
\end{table*}

\paragraph{Operational interpretation.}
Retrieval names the function of obtaining sources; in this implementation the model can generate them. We do not observe retrieval from internal memory. Mapping aligns the retrieved \emph{sources with each other}. It does not receive the target. Schema induction receives the mapping rather than the raw stories, making explicit whether the intermediate representation is sufficient. Adaptation receives the target again, but not the raw sources or mapping. These restrictions create a traceable source-comparison-to-schema-to-target path. They are experimental choices, not a universal ordering of human analogy.

For example, sources implementing convergence should preserve the constraints that a concentrated force is unsafe along a path, individual weak forces are insufficient, and coordinated forces combine at the goal. Mapping should retain those relations; schema induction should state their conditions without relying on story-specific names; adaptation should check whether the target permits the required combination. ``Organize resources better'' would be too general to qualify as a useful schema.

\paragraph{Diagnosis and cumulative repair.}
We inspect naturally unsuccessful trajectories in stage order and record the first output inadequate for its assigned function. A reference-informed author, GPT-6 Sol, constructs a corrected output for that stage while earlier outputs are retained. The author sees the target reference answer, an \emph{open-oracle} construction setting. Upstream corrections must be defensible under the stage's information contract and must not explicitly reveal the target answer. The same base solver then reruns only the downstream stages without seeing the reference answer.

If the target is solved, the continuation has been rescued under that correction. Otherwise, we inspect the new trajectory, retain previous corrections, and correct the next inadequate stage. This cumulative procedure reveals failures hidden behind upstream deficiencies. Rescue is conditional on the repaired path and does not estimate the isolated causal effect of each stage across all problems. A gold adaptation output already contains the final answer and has no downstream continuation; its reference agreement is reported separately from base-model rescue.

We report Llama-3.1-8B-Instruct on 100 GPQA Main items and 100 MuSR murder mysteries. Human review covers 30\% of questions and associated outputs for logical coherence and correctness. The remaining cases lack human review; oracle-informed construction is not equivalent to exhaustive validation. Diagnosis concerns observed artifacts in the imposed workflow, rather than hidden reasoning mechanisms.
